\exercisesection{}

\begin{enumerate}
\def\labelenumi{\arabic{enumi})}
\item
  设 E 为 R、C 或 H 上的有限维向量空间，\(\Phi\) 和 \(\Phi'\) 为 E 上两个非退化正厄米形式。假设 \(\mathbf{U}(\Phi) \subset \mathbf{U}(\Phi')\)。证明 \(\Phi\) 与 \(\Phi'\) 成比例。（利用如下事实：存在一个标准正交基 \((e_{1}, \ldots, e_{n})\)，它关于 \(\Phi\) 是正交的，并且关于 \(\Phi'\) 也是正交的。映射 \(u \in \mathcal{L}(\mathrm{E}, \mathrm{E})\) 满足 \(u(e_{i}) = e_{j}\)、\(u(e_{j}) = e_{i}\)、\(u(e_{k}) = e_{k}\)（对 \(k \neq i, j\)），且属于 \(\mathbf{U}(\Phi)\)。）
\item
  采用引理 7 的记号。证明对于 Haar 测度，\(\Omega\) 在 \(\mathbf{GL}(n,\mathbf{K})\) 中的补集可忽略。（仿照命题 6 论证。）
\end{enumerate}

¶ 3) 设 X 为局部紧空间，局部紧群 H 通过 \((x,\xi)\mapsto x\xi\) 在其中连续且适当地右作用（\(x\in X\)，\(\xi\in H\)）。令 \(\pi\) 为典范映射 \(X\to X/H\)。令 \(\rho\) 为 H 在 \(\mathbf{GL}(n,\mathbf{C})\) 中的连续表示。证明对每个 \(b\in X/H\)，都存在 X/H 中 b 的一个邻域 U 以及从 \(\overline{\pi}(\mathrm{U})\) 到 \(\mathbf{GL}(n,\mathbf{C})\) 的连续映射 r，使得 \(r(x\xi)=r(x)\rho(\xi)\) 对所有 \(x\in\overline{\pi}(\mathrm{U})\) 和 \(\xi\in H\) 成立。（可以假设 X/H 是紧的。令 f 为 X 上满足 \(\geqslant0\)、在任何轨道上都不恒为零且具有紧支集的连续函数。置

\[
r (x) = \int_ {\mathbf {H}} f (x \xi) \rho (\xi) ^ {- 1} d \beta (\xi) \in \mathbf {M} _ {n} (\mathbf {C}),
\]

其中 \(\beta\) 表示 H 上的左 Haar 测度。证明适当选取 \(f\) 和 U 后，\(r(x) \in \mathbf{GL}(n, \mathbf{C})\) 对每个 \(x \in \overline{\pi}^1(\mathrm{U})\) 成立。

\begin{enumerate}
\def\labelenumi{\arabic{enumi})}
\setcounter{enumi}{3}
\tightlist
\item
  设 \(\mathbf{K}\) 是一个非离散交换局部紧域。设 \(\mathbf{A}\) 是 \(\mathbf{K}\) 上的有限秩代数。
\end{enumerate}

\begin{enumerate}
\def\labelenumi{\alph{enumi})}
\tightlist
\item
  令 \(\mathrm{T}(n, \mathbf{A})\) 为由矩阵 \((x_{ij}) \in \mathbf{M}_n(\mathbf{A})\) 构成的代数，其中 \(x_{ij} = 0\) 当 \(i > j\) 时。证明若 \(M = (m_{ij}) \in \mathrm{T}(n, \mathbf{A})\)，则
\end{enumerate}

\[
\mathrm{N} _ {\mathrm{T} (n, \mathrm{A}) / \mathrm{K}} (M) = \prod_ {i = 1} ^ {n} \mathrm{N} _ {\mathrm{A} / \mathrm{K}} (m _ {i i} ^ {i - n + 1}).
\]

（使用引理 6。）

\begin{enumerate}
\def\labelenumi{\alph{enumi})}
\setcounter{enumi}{1}
\item
  令 \(\mathrm{T}(n,\mathrm{A})^*\) 为 \(\mathrm{T}(n,\mathrm{A})\) 的可逆元素群。设 \(M = (m_{ij}) \in \mathrm{T}(n,\mathrm{A})\)。证明 \(M \in \mathrm{T}(n,\mathrm{A})^*\) 当且仅当 \(m_{ii}\) 在 \(\mathbf{A}\) 中可逆（对所有 \(i\)）。
\item
  证明 \(\mathbf{T}(n,\mathbf{A})^{*}\) 上的左 Haar 测度由下式给出：
\end{enumerate}

\[
\left(\prod_ {i = 1} ^ {n} \operatorname{mod} \mathrm{N} _ {\mathrm{A} / \mathrm{K}} (m _ {i i}) ^ {i - n - 1}\right) \cdot \bigotimes_ {i \leqslant j} d \alpha (m _ {i j}),
\]

其中 \(\alpha\) 表示 A 的加法群上的 Haar 测度。

\begin{enumerate}
\def\labelenumi{\alph{enumi})}
\setcounter{enumi}{3}
\tightlist
\item
  证明若 \(\mathbf{N}_{\mathbf{A} / \mathbf{K}} = \mathbf{N}_{\mathbf{A}^0 /\mathbf{K}}\)，则
\end{enumerate}

\[
\Delta_ {\mathrm{T} (n, \mathrm{A}) ^ {*}} \left(\left(m _ {i j}\right)\right) = \prod_ {i = 1} ^ {n} \operatorname{mod} \mathrm{N} _ {\mathrm{A} / \mathrm{K}} \left(m _ {i i}\right) ^ {2 i - n - 1}.
\]

\begin{enumerate}
\def\labelenumi{\arabic{enumi})}
\setcounter{enumi}{4}
\tightlist
\item
  在 \(\mathbf{R}^{n}\) 上赋予二次型 \(x_{1}^{2} + x_{2}^{2} + \cdots + x_{n}^{2}\)。令 \(G_{n}\) 为行列式为 1 的 \(\mathbf{R}^{n}\) 的位移群（Alg., Ch. IX, §6, No.~6, Def. 3）。
\end{enumerate}

\begin{enumerate}
\def\labelenumi{\alph{enumi})}
\item
  证明 \(\mathbf{G}_n\) 是幺模的。
\item
  \(\mathbf{G}_2\) 的每个元素都能唯一写成如下形式：
\end{enumerate}

\[
(x, y) \mapsto (u + x \cos \omega - y \sin \omega , v + x \sin \omega + y \cos \omega),
\]

其中 \(u, v\) 属于 \(\mathbf{R}\)，而 \(\omega\) 是半直线角群 \(\Theta\) 的元素，且该群同构于 \(\mathbf{U}\)。证明 \(du \otimes dv \otimes d\omega\)（其中 \(d\omega\) 是 \(\Theta\) 上的 Haar 测度）是 \(G_2\) 上的 Haar 测度。

\begin{enumerate}
\def\labelenumi{\arabic{enumi})}
\setcounter{enumi}{5}
\tightlist
\item
  令 \(\mathbf{P}\) 为复数 \(z = x + iy\) 所成的集合，其虚部为 \(>0\)。证明 \(\mathbf{SL}(2,\mathbf{R})\) 按下式在 \(\mathbf{P}\) 上连续且传递地左作用：
\end{enumerate}

\[
\left( \begin{array}{c c} a & b \\ c & d \end{array} \right) \cdot z \mapsto \frac {a z + b}{c z + d}
\]

并证明由此 P 成为 \(\mathbf{SL}(2,\mathbf{R})\) 的拓扑齐性空间。证明 \(y^{-2}dx dy\) 是 P 上的不变测度。

¶ 7) 设 G 为幺模群 \(\mathbf{SL}(n,\mathbf{R})\)，\(\Gamma\) 为由行列式为 1 且元素为整数的矩阵构成的 G 的子群。证明通过对 n 作归纳，齐性空间上每个不变测度 \(\mu\) 在 \(G/\Gamma\) 上都是有界的，并且对 \(f \in \mathcal{K}(\mathbf{R}^{n})\)，适当选取 \(\mu\) 后有

\[
\int_ {\mathbf {R} ^ {n}} f (x) d x = \int_ {\mathrm{G} / \Gamma} \left(\sum_ {z \in \mathbf {Z} ^ {n}, z \neq 0} f (X \cdot z)\right) d \mu (\dot {X})\tag{1}
\]

\((\dot{X} = X\Gamma)\)。令 \(G'\) 为 \(G\) 的子群，保持基向量 \((1,0,0,\ldots,0)\) 不变，从而可将 \(G/G'\) 与 \(\mathbf{R}^n\) 识别。令 \(G'' = G' \cap \Gamma\)。令 \(H\) 为由形如 \(\begin{pmatrix} 1 & w \\ 0 & Y \end{pmatrix}\) 的 \(G'\) 中的矩阵组成的子群，其中 \(Y \in \mathbf{SL}(n-1, \mathbf{R})\) 的元素为整数。则 \(H/G''\) 是紧的，并且根据归纳假设，\(G'/H\) 具有有限不变测度，因此根据命题 12, §2, No.~8，\(G'/G''\) 具有有限不变测度。应用 §2 的 Exer. 6 以及 A, VII, §4 的 Exer. 20 b)，证明对 \(f \in \mathcal{K}(\mathbf{R}^n)\)，

\[
a \int_ {\mathbf {R} ^ {n}} f (x) d x = \int_ {\mathrm{G} / \Gamma} \left(\sum f (X \cdot m)\right) d \mu (\dot {X}),
\]

其中求和遍及向量 \(m = (m_1, \ldots, m_n)\) 的集合，这些向量的坐标是互素的整数（A, VII, §1, No.~2, Th. 1）。将此应用于函数 \(f(2x), f(3x), \ldots\) 并求和，证明 (1)。然后将 \(f\) 替换为函数 \(x \mapsto \varepsilon^n f(\varepsilon x)\)，其中 \(f \geqslant 0\)，并令 \(\varepsilon\) 趋于 0，证明 \(G/\Gamma\) 的每个紧子集的测度都 \(\leqslant 1\)。

\begin{enumerate}
\def\labelenumi{\arabic{enumi})}
\setcounter{enumi}{7}
\tightlist
\item
  \begin{enumerate}
  \def\labelenumii{\alph{enumii})}
  \tightlist
  \item
    证明在 \(\mathbf{S}_{n-1}\) 上存在一个测度 \(\omega_{n-1}\)，它在 \(\mathbf{O}(n,\mathbf{R})\) 下不变，且除相差一个常数因子外只有一个（将 \(\mathbf{S}_{n-1}\) 看作 \(\mathbf{O}(n,\mathbf{R})\) 的齐性空间）。b) 映射 \((t,z)\mapsto tz\) 允许将拓扑空间 \(\mathbf{R}^n -\{0\}\) 与拓扑空间 \(\mathbf{R}_+^*\times \mathbf{S}_{n-1}\) 识别。证明在 \(\mathbf{R}^n -\{0\}\) 上由 \(\mathbf{R}^n\) 上的 Lebesgue 测度诱导的测度，可以在相差一个常数因子的意义下与 \(t^{n-1}dt\otimes d\omega_{n-1}(z)\) 识别。（利用 \(\mathbf{R}^n\) 上的 Lebesgue 测度在 \(\mathbf{O}(n,\mathbf{R})\) 下不变这一事实，以及中心为 0、半径为 \(r\) 的闭球的 Lebesgue 测度与 \(r^n\) 成比例这一事实。）
  \end{enumerate}
\end{enumerate}

\begin{enumerate}
\def\labelenumi{\alph{enumi})}
\setcounter{enumi}{2}
\item
  令 \(L_{h}\) 为 \(S_{n-1}\) 与超平面 \(x_{n}=h\) 的交。非空的 \(L_{h}\) 是 \(\mathbf{O}(n-1,\mathbf{R})\) 在 \(S_{n-1}\) 中的不变类。证明 \(\omega_{n-1}\) 可以写成 \(\int\lambda_{h}d\nu(h)\) 的形式，其中 \(\lambda_{h}\) 是 \(\mathbf{O}(n-1,\mathbf{R})\) 下不变且支承于 \(L_{h}\) 的测度，而 \(\nu\) 是 \([-1,1]\) 上的测度。将 \(L_{h}\) 与 \(S_{n-2}\) 通过 \(z\mapsto he_{n}+\sqrt{1-h^{2}}z\) 识别，可设 \(\lambda_{h}=\omega_{n-2}\)。令 \(K_{\theta}\) 为 \(S_{n-1}\) 的由 \(\sin\theta\leqslant x_{n}\leqslant1\) 定义的子集。通过计算 tz 的集合（\(0\leqslant t\leqslant1\)，\(z\in K_{\theta}\)）的 Lebesgue 测度，证明 \(\omega_{n-1}(K_{\theta})\) 与 \(\int_{\theta}^{\pi/2}\cos^{n-2}\varphi d\varphi\) 成比例。由此推出，若令 \(h=\sin\theta(-\pi/2\leqslant\theta\leqslant\pi/2)\)，则测度 \(\nu\) 可以与 \(\cos^{n-2}\theta d\theta\) 识别。
\item
  由对 \(n\) 的归纳，推出
\end{enumerate}

\[
d \omega_ {n - 1} = \cos^ {n - 2} \theta_ {1} \cos^ {n - 3} \theta_ {2} \dots \cos \theta_ {n - 2} d \theta_ {1} d \theta_ {2} \dots d \theta_ {n - 1},
\]

 其中

\[
\begin{array}{r l} & x _ {1} = \sin \theta_ {1} \\ & x _ {2} = \cos \theta_ {1} \sin \theta_ {2} \\ & \dots \\ & x _ {n - 1} = \cos \theta_ {1} \cos \theta_ {2} \dots \cos \theta_ {n - 2} \sin \theta_ {n - 1} \\ & x _ {n} = \cos \theta_ {1} \cos \theta_ {2} \dots \cos \theta_ {n - 2} \cos \theta_ {n - 1} \end{array}
\]

其中 \(-\pi / 2 \leqslant \theta_{i} \leqslant \pi / 2\) 适用于 \(1 \leqslant i \leqslant n - 2, 0 \leqslant \theta_{n-1} < 2\pi\)。

\begin{enumerate}
\def\labelenumi{\arabic{enumi})}
\setcounter{enumi}{8}
\tightlist
\item
  令 \((e_1, e_2, e_3)\) 为 \(\mathbf{R}^3\) 的标准基。每个矩阵 \(\sigma \in \mathbf{SO}(3, \mathbf{R})\)，只要满足 \(\sigma(e_3) \neq e_3\) 且 \(\sigma(e_3) \neq -e_3\)，都能唯一写成
\end{enumerate}

\[
\begin{array}{l} \sigma (\varphi , \psi , \theta) = \\ \left( \begin{array}{c c c} \cos \varphi   \cos \psi - \sin \varphi   \sin \psi   \cos \theta & - \cos \varphi   \sin \psi - \sin \varphi   \cos \psi   \cos \theta & \sin \varphi   \sin \theta \\ \sin \varphi   \cos \psi + \cos \varphi   \sin \psi   \cos \theta & - \sin \varphi   \sin \psi + \cos \varphi   \cos \psi   \cos \theta & - \cos \varphi   \sin \theta \\ \sin \psi   \sin \theta & \cos \psi   \sin \theta & \cos \theta \end{array} \right), \end{array}
\]

其中 \(0 < \theta < \pi\)、\(0 \leqslant \varphi < 2\pi\)、\(0 \leqslant \psi < 2\pi\)（欧拉角）。证明 \(\sin \theta d\theta d\varphi d\psi\) 是 \(\mathbf{SO}(3, \mathbf{R})\) 上的 Haar 测度。（将齐性空间 \(\mathbf{SO}(3, \mathbf{R}) / \mathbf{SO}(2, \mathbf{R})\) 与 \(\mathbf{S}_2\) 识别，使用 Exer. 8 以及 §2 的 Th. 3。）

\begin{enumerate}
\def\labelenumi{\arabic{enumi})}
\setcounter{enumi}{9}
\tightlist
\item
  令 D 为行列式为 1 的 \(\mathbf{R}^{3}\) 的位移群，即 \(\mathbf{SO}(3,\mathbf{R})\) 与平移群 T 的半直积。采用 Exer. 9 的记号 \(\sigma(\varphi,\psi,\theta)\)，并以 \(t(\xi,\eta,\zeta)\) 表示平移向量为 \((\xi,\eta,\zeta)\) 的平移。
\end{enumerate}

\begin{enumerate}
\def\labelenumi{\alph{enumi})}
\item
  令 H 为 D 中保持 \(\mathbf{R}e_{3}\) 稳定并保持 \(\mathbf{R}e_{3}\) 上定向的闭子群。证明 H 是平移群 \(t(0,0,\zeta)\) 与旋转群 \(\sigma(\omega,0,0)\) 的乘积。由此推出，齐性空间 E = D/H（可将其与 \(\mathbf{R}^{3}\) 的定向仿射直线空间识别）具有一个在 D 下不变的测度，且除相差一个常数因子外只有一个。（使用 Exer. 5 a)。
\item
  令 \(\mathbf{E}_1\) 为 \(\mathbf{E}\) 中由不垂直于 \(e_3\) 的定向直线构成的开子空间；这样的定向直线可由其交点的坐标 \(\xi', \eta'\)（交点位于 \(\mathbf{R}e_{1} + \mathbf{R}e_{2}\)），以及方向向量的坐标 \((\sin \varphi \sin \theta, -\cos \varphi \sin \theta, \cos \theta)\) 确定。证明 \(E_{1}\) 上在 D 下不变的测度由下式给出：
\end{enumerate}

\[
\sin \theta | \cos \theta | d \xi^ {\prime} d \eta^ {\prime} d \varphi d \theta .
\]

\begin{enumerate}
\def\labelenumi{\arabic{enumi})}
\setcounter{enumi}{10}
\tightlist
\item
  设 \(G\) 为紧群 \(\mathbf{O}(n, \mathbf{R})\)，\(\lambda\) 为 \(G\) 上的归一化 Haar 测度。对 \(0 < k < n\)，令 \(H\) 为 \(G\) 中的闭子群，它（全局）保持 \(\mathbf{R}^k\)（\(\mathbf{R}^n\) 中）不变；赋予商拓扑的齐性空间 \(G / H\) 可典范地与 Grassmann 流形 \(E = G_{n-1, k-1}(\mathbf{R})\) 识别，该流形由 \(k\) 维的 \(\mathbf{R}^n\) 线性子空间组成。存在一个测度 \(\mu\)，定义在 \(E\) 上且在 \(G\) 下不变，除相差一个常数外确定。对每个子空间 \(P \in E\)，令 \(\sigma_P\) 为 \(\mathbf{R}^k\) 的 Lebesgue 测度在 \(s \in G\) 下的像，其中 \(s \cdot \mathbf{R}^k = P\)（与满足此关系的 \(s \in G\) 的选取无关）。证明可以如此选取 \(\mu\)，使其满足如下性质：对每个连续函数 \(f\)，其定义在 \(\mathbf{R}^n\) 上且具有紧支集，以及每个 \(P \in E\)，令 \(F(P) = \int_{P} f(x) d\sigma(x)\)；则
\end{enumerate}

\[
\int_ {\mathrm{E}} \mathrm{F} (\mathrm{P}) d \mu (\mathrm{P}) = \int_ {\mathbf {R} ^ {n}} \| x \| ^ {k - n} f (x) d x.
\]

（若 \(\mathbf{P}_0 = \mathbf{R}^k\)，注意可以写成 \(\int_{\mathbf{E}} \mathrm{F}(\mathrm{P}) d\mu(\mathrm{P}) = c \int_{\mathbf{G}} \mathrm{F}(s \cdot \mathrm{P}_0) d\lambda(s)\)，其中 \(c\) 是适当的常数；另一方面，

\[
\int_ {\mathrm{G}} f (s \cdot x) d \lambda (s) = c ^ {\prime} \int_ {\mathbf {S} _ {n - 1}} f (\| x \| z) d \omega_ {n - 1} (z)
\]

采用习题 8 的记号；最后使用 Exer. 8 b)。）

¶ 12) 设 K 为交换域，E 为 K 上的向量空间，F 为 E 的线性子空间，p 为 E 到 E/F 的典范同态，A 为满足 \(f: E/F \to E\) 且 \(p \circ f\) 是 E/F 的恒等同态的同态所成的集合。

\begin{enumerate}
\def\labelenumi{\alph{enumi})}
\item
  若 \(f \in A\) 且 \(h \in \operatorname{Hom}(\mathrm{E} / \mathrm{F}, \mathrm{F})\)，则 \(f + h \in A\)。证明若 \(f, f' \in A\)，则存在唯一的 \(h \in \operatorname{Hom}(\mathrm{E} / \mathrm{F}, \mathrm{F})\)，使得 \(f + h = f'\)。因此可将 \(A\) 看作一个仿射空间，其平移空间为 \(\operatorname{Hom}(\mathrm{E} / \mathrm{F}, \mathrm{F})\)。
\item
  令 B 为 E 中与 F 互为补空间的线性子空间所成的集合。证明 \(f \mapsto f(\mathrm{E}/\mathrm{F})\) 是从 A 到 B 的双射。因此可将 B 看作一个仿射空间，其平移空间为 \(\operatorname{Hom}(\mathrm{E}/\mathrm{F}, \mathrm{F})\)。
\item
  证明若 \(u\) 是 \(E\) 的自同构且 \(u(F) = F\)，则映射 \(f \mapsto u \circ f\) 是从 \(A\) 到 \(A\) 的仿射双射。（在 \(A\) 中选取一个原点，并注意映射 \(h \mapsto u \circ h\) 是向量空间 \(\operatorname{Hom}(E/F, F)\) 的自同构。）由此推出映射 \(F' \mapsto u(F')\) 是从 \(B\) 到 \(B\) 的仿射双射。
\item
  假设 K = R 且 \(\dim E < +\infty\)。令 G 为紧群，\(\rho\) 为 G 在 E 中的连续线性表示，满足 \(\rho(s)(F) = F\) 对所有 \(s \in G\)。证明存在 \(F' \in B\)，使得 \(\rho(s)(F') = F'\) 对所有 \(s \in G\)。（使用 c) 和 No.~2 的引理 2。）再利用 No.~1 的命题 1 重新得到此结果。
\end{enumerate}

